\section{The construction of SL(2,R)/SO(2) coset}

We briefly review the construction of the $\frac{SL(2, R)}{SO(2)}$ coset representative, which is important for discussion of the SL(2, R) symmetry, and the interplay between SL(2, Z) and U(1)-violation.

The coset representative of $\frac{SL(2, R)}{SO(2)}$ is an SL(2, R) matrix $\mathcal{V}$ with a local SO(2) symmetry. 
Generically, $\dim SL(2, R) = 3$ so $\mathcal{V}$ is parameterised by three real scalars: $\mathcal{V} = \mathcal{V}(\tau_1, \tau_2, \theta)$.
It is convenient to separate the SO(2) factor, parameterised by $\theta$, using the Iwasawa decomposition \cite{iwasawa1949}
\begin{equation}
\mathcal{V} = 
\begin{pmatrix}
1 & 0 \\
\tau_1 & 1
\end{pmatrix}
\begin{pmatrix}
\frac{1}{\sqrt{\tau_2}} & 0\\
0 & \sqrt{\tau_2}
\end{pmatrix}
\begin{pmatrix}
\cos\theta & \sin\theta\\
-\sin\theta & \cos\theta
\end{pmatrix}.
\label{Iwasawa-decomposition}
\end{equation}
In this representation, SL(2, R) acts on $\mathcal{V}$ from the left, with $\Lambda \equiv \begin{pmatrix}
d&c\\b&a\end{pmatrix} \in SL(2, R)$.
The representative $\mathcal{V}$ by itself is not physical due to the SO(2) gauge redundancy.
However, one can construct $\mathcal{V}\mathcal{V}^T$ which transforms covariantly under SL(2, R) and removes the SO(2) factor and gives the physical metric for the scalar manifold, parameterised by $\tau_1, \tau_2$:
\begin{equation}
\mathcal{M} = \mathcal{V}\mathcal{V}^T = \frac{1}{\tau_2}\begin{pmatrix}
1 & \tau_1\\
\tau_1 & |\tau|^2
\end{pmatrix}.
\end{equation}


To construct the kinetic term for the scalars, one builds the sl(2, R)-valued Maurer-Cartan 1-form $\mathcal{V}^{-1}\partial_m\mathcal{V}$, which decomposes into the non-compact and compact elements of the sl(2, R) algebra: $\mathcal{V}^{-1}\partial_m\mathcal{V} = P_m + Q_m$, where $P_m$ projects onto the non-compact generators and $Q_m$ is the SO(2) connection.
The kinetic term for the scalars is then given by $tr P^m P_m$, with positive definite as ensured by tracing over non-compact elements of a Lie algebra.
Computing $P, Q$, with $SO(2)\cong U(1)$ compexified via $e^{i\theta} = \cos\theta + i\sin\theta$, one finds
\begin{equation}
P = P_1\begin{pmatrix}
    0&1\\1&0
\end{pmatrix}+P_2\begin{pmatrix}
    1&0\\0&-1
\end{pmatrix},\quad
P_1 + iP_2 = e^{2i\theta}\frac{\partial \tau}{2 \tau_2},\quad
Q = \partial \theta - \frac{1}{2}\frac{\partial \tau_1}{\tau_2}.
\end{equation}
So the vielbein $P$ transforms covariantly local U(1) with charge 2, and that the kinetic term $\frac{|\partial \tau|^2}{2\tau_2^2}$ is U(1) invariant.
This U(1) is independent of the global SL(2, R) transformation, and is a gauge redundancy reflecting the choice of the U(1) frame.

The parameter $\theta$ here is the non-dynamical U(1) phase. 
We can remove its presence in the action by picking a gauge $\theta = 0$, leaving two real scalars $\tau_1, \tau_2$ that parameterise the 2d coset $\frac{SL(2, R)}{U(1)}$ without redundancy.
However, general SL(2, R) transformations do not preserve the $\theta = 0$ gauge, hence a compensating U(1) transformation given by $\theta = - \arg{(c\tau+d)}$ is induced when  acts on $\mathcal{V}$ from the left.
This is independent of the local U(1) symmetry which may be violated, but part of the global SL(2, Z) symmetry which should always be preserved.

This is another way to view the breaking of the U(1) symmetry: the original U(1) being the SO(2) factor in the SL(2, R), hence when SL(2, R) is broken to SL(2, Z), the U(1) is also broken.

% From the SL(2, R) representative $\mathcal{V}$, one can construct an U(1)-invariant, symmetric metric on the scalar manifold $\mathcal{M} = \mathcal{V} \mathcal{V}^T  = \frac{1}{\tau_2}\begin{pmatrix}
% |\tau|^2 & \tau_1\\
% \tau_1 & 1
% \end{pmatrix}$, which transforms under SL(2, R) as $\mathcal{M} \to \Lambda \mathcal{M} \Lambda^T$, for $\Lambda = \begin{pmatrix}
% a&b\\c&d
% \end{pmatrix} \in SL(2, R)$.

Besides $P$, the other bosonic field that transforms under U(1) is the 3-form field strength $G_3 = e^{\Phi/2}(F_3 - \tau H_3)$, which carries charge 1 under U(1).
When computing higher-derivative corrections, one wishes to keep track of the U(1) charge for possible U(1)-violation corrections, in that context it is convenient to work with the complexified vielbein $P$ and 3-form field strength $G_3$.
It is important to distinguish this U(1) from the compensating U(1) to remove the redundant scalar $\theta$ in the coset representative $\mathcal{V}$. The compensating U(1) is an SL(2, Z) transformation, which the effective action should always be invariant under. For example, the modular form $f(\tau, \bar\tau)$ also carries a U(1)-like charge under SL(2, Z) but it does not carry local U(1) charge, and hence may contribute to SL(2, Z) invariant but U(1)-violating corrections.